\chapter{Alloy Model Details} \label{app:alloy_details}

This appendix provides the model architecture, per-module descriptions, and
detailed per-covenant verification results summarized in
Chapter~\ref{ch:verification}.

% ======================================================================
\section{Model Architecture}

The Alloy models follow a three-layer inheritance hierarchy
(Figure~\ref{fig:alloy_hierarchy}).

\begin{figure}[t]
\centering
\begin{tikzpicture}[
  box/.style={rectangle, draw, rounded corners=3pt, minimum width=2cm,
              minimum height=0.7cm, font=\footnotesize, align=center},
  base/.style={box, fill=gray!15},
  concrete/.style={box, fill=blue!8},
  cross/.style={box, fill=orange!10, dashed},
  >=stealth, thick,
  node distance=0.6cm,
]
  \node[base] (btc) {\texttt{btc\_base}};
  \node[base, right=2cm of btc] (vault) {\texttt{vault\_base}};
  \draw[->] (btc) -- node[above, font=\scriptsize] {opens} (vault);

  \node[concrete, below=1.8cm of vault, xshift=-4cm] (ctv)
    {\texttt{ctv\_vault}};
  \node[concrete, right=0.25cm of ctv] (ccv)
    {\texttt{ccv\_vault}};
  \node[concrete, right=0.25cm of ccv] (opv)
    {\texttt{opvault}};
  \node[concrete, right=0.25cm of opv] (cat)
    {\texttt{cat\_csfs}};
  \node[concrete, right=0.25cm of cat] (sim)
    {\texttt{simplicity}};

  \foreach \c in {ctv, ccv, opv, cat, sim}
    \draw[->] (vault) -- (\c);

  \node[cross, right=1cm of sim, yshift=0.9cm] (threat)
    {\texttt{threat\_model}};
  \draw[->, dashed] (threat) -- (sim);
  \draw[->, dashed] (threat.south west) -- (cat.north east);

  \node[cross, below=1cm of opv] (cross)
    {\texttt{cross\_covenant}};
  \draw[->, dashed] (ccv) -- (cross);
  \draw[->, dashed] (opv) -- (cross);

  \node[font=\scriptsize, text=gray, anchor=west] at (threat.east)
    [xshift=0.2cm] {\itshape used by all};
\end{tikzpicture}
\caption{Alloy model hierarchy. \texttt{btc\_base} defines UTXOs, keys, and
transactions. \texttt{vault\_base} defines the abstract vault state machine.
Five concrete models extend the base with covenant-specific transition guards.
\texttt{threat\_model} and \texttt{cross\_covenant} are cross-cutting modules.}
\label{fig:alloy_hierarchy}
\end{figure}

\paragraph{\texttt{btc\_base.als}.}
Defines the Bitcoin transaction model: signatures for \texttt{UTXO}, \texttt{Key},
\texttt{Transaction}, and \texttt{Address}. UTXOs have an amount, owner key, and
spent status. Transactions consume inputs and produce outputs. This module provides
the vocabulary shared by all vault models.

\paragraph{\texttt{vault\_base.als}.}
Defines the abstract vault state machine: four states (\texttt{VAULTED},
\texttt{UNVAULTING}, \texttt{WITHDRAWN}, \texttt{RECOVERED}), transitions between
states (trigger, withdraw, recover), and the \texttt{VaultFamily} concept that groups
UTXOs belonging to the same vault instance. The CSV delay is modeled as an integer
field on each vault UTXO. Properties defined at this level (fund conservation,
recovery liveness) apply uniformly to all concrete models.

\paragraph{Concrete vault models.}
Each of the five models extends \texttt{vault\_base} with covenant-specific
transition guards. \texttt{ccv\_vault} models the CCV mode enumeration
and the OP\_SUCCESS behavior for undefined modes. \texttt{opvault\_vault} models
the three-key separation and authorized recovery. \texttt{cat\_csfs\_vault} models
the dual-verification binding and unconstrained recovery leaf.

\paragraph{\texttt{threat\_model.als}.}
Defines attacker profiles by specifying which keys the attacker controls. Properties
are checked under each profile to determine the minimum capability needed for each
attack class.

\paragraph{\texttt{cross\_covenant.als}.}
Models interactions between vault instances from different covenant types sharing the
same UTXO set. Checks whether cross-vault injection is possible and whether fee
wallet contention between concurrent OP\_VAULT operations can block recovery.

% ======================================================================
\section{Per-Covenant Detailed Results}

\paragraph{CTV.}
All six checks pass: no unauthorized extraction (with no keys and with hot key only),
CSV enforcement, no state proliferation, eventual withdrawal, and recovery liveness.
The address reuse scenario finds an instance where a second deposit creates an
unspendable UTXO, confirming the structural vulnerability (TM7).

\paragraph{CCV.}
All seven checks pass, including the mode confusion containment check (undefined
modes trigger OP\_SUCCESS but only within the affected leaf---the vault's other
leaves remain secure). The griefing scenario demonstrates that any party can invoke
recovery without keys (TM2). The splitting scenario shows the revault chain that
enables the per-UTXO recovery-cost scaling of class Z3 (TM4).

\paragraph{OP\_VAULT.}
All eight checks pass, including the dual-key-no-theft property (trigger +
recoveryauth compromise produces liveness denial but not fund theft, because the
recovery address is pre-committed). The key-loss scenario demonstrates permanent
inability to recover if the recoveryauth key is destroyed.

\paragraph{CAT+CSFS.}
Five of six checks pass. The \texttt{catcsfsNoExtraction\_ColdKeyOnly} check
\emph{fails}: the Alloy Analyzer finds a trace where the cold key holder invokes
recovery with an attacker-controlled output address, extracting funds. This is the
structural confirmation of TM11 and the only expected assertion failure across all
models. The destination lock and CSV enforcement checks pass.

\paragraph{Simplicity.}
All six checks pass: no extraction under any single-key compromise, output binding
enforcement, CSV enforcement, and no state proliferation. The recovery-constrained
scenario confirms that cold key compromise can only recover to the pre-committed
address (unlike CAT+CSFS).

% ======================================================================
\section{Limitations of the Formal Models}

The Alloy models operate at the state-machine level, abstracting away cryptographic
primitives, serialization details, and temporal dynamics. Signatures are modeled as
key-possession checks (if the attacker holds the key, the signature succeeds). Hash
functions are assumed collision-resistant. Network propagation, mining centralization,
and mempool dynamics are not modeled. These abstractions mean the Alloy verification
confirms \emph{structural} properties (state reachability, authorization requirements)
but not \emph{cryptographic} properties (signature unforgeability, hash binding) or
\emph{temporal} properties (race conditions, confirmation timing). The regtest
experiments partially compensate for the temporal gap by demonstrating actual
transaction construction and broadcast, while the cryptographic properties are
inherited from the underlying BIP-340 Schnorr and SHA-256 specifications.
